\chapter*{Related Work}
\addcontentsline{toc}{chapter}{Related Work}

AAT is a mathematical theory that constructs relative architecture from Atoms and Laws.
It defines objects and structure-preserving morphisms relative to a reading, and develops
their geometry, transport, comparison, and local reconstruction. Below we compare objects,
assumptions, morphisms, and conclusions, focusing on algebraic specification, sheaf
semantics, abstract interpretation, data migration, and bidirectional transformation, and
situate the constructions and consequences proper to AAT.

\section{Semantics of specifications and architecture}\label{sec:R.1}

When the description of a specification changes, its meaning must change accordingly.
The institutions of Goguen and Burstall define this relation independently of the logic.
An institution carries signatures, sentences, categories of models, and a satisfaction
relation, and requires that the translation of sentences and the reduction of models along
a change of signature preserve the satisfaction relation. Morphisms between models are also
part of this framework
(\cite[\href{https://courses.grainger.illinois.edu/cs522/sp2016/InstitutionsAbstractModelTheory.pdf}{\S2.1, Definition 1}]{GB92}).
Moreover, colimits of theories can be constructed from colimits of signatures.
This result makes it possible to state the principle of assembling large specifications from
small ones on a foundation common to different logics
(\cite[\href{https://courses.grainger.illinois.edu/cs522/sp2016/InstitutionsAbstractModelTheory.pdf}{Theorem 11 and its consequences, p. 108}]{GB92}).

Once the connection of parts is taken as the object of study, compatibility of behavior
becomes the issue. Allen and Garlan's Wright describes the ports of components and the roles
and glue of connectors in CSP. If a connector is conservative and deadlock-free, and
the ports attached to each role satisfy the compatibility condition, then the connector
instantiated with those ports is deadlock-free as well
(\cite[\href{https://ix.cs.uoregon.edu/~michal/Classes/f01/fsv/papers/Allen-Garlan.pdf}{\S\S8.1--8.2, Definition 8.1.2, Theorem 8.2.3}]{AG97}).

AAT fixes a reading of Atoms and determines the relative core from the values of the Laws.
Adding covers, coefficients, and evaluation, it constructs a geometry and, through equations
and ideals, expresses the satisfaction of Laws as a geometric zero condition and defines the
lawful locus (Theorem~\ref{thm:2.9}).
This geometry is the object on which obstructions to repair, transport along readings, and
comparison of construction routes are treated.

For protocols, that correspondence can be captured by preservation equations for operations
and observations. Once named operations, the relations among paths, and the observations are
fixed, a morphism between realizations is a vertexwise state map that preserves the
operations and the observations. AAT constructs a typed input from this independent
semantics and puts the morphisms on the two sides in correspondence
(Propositions~\ref{prop:1.38} and~\ref{prop:1.43}).

\section{Local consistency and cohomological obstructions}\label{sec:R.2}

Behaviors given part by part are read as the behavior of a single system.
Goguen's sheaf semantics is a precedent for this local-to-global description.
It treats objects as sheaves and interactions as diagrams, and, under suitable completeness,
obtains the behavior of a system as the limit of a diagram
(\cite[\href{https://ncatlab.org/nlab/files/Goguen-SheafSemantics.pdf}{public manuscript, \S\S2, 3.1--3.3, Proposition 26}]{Goguen92}).
Gibson's preprint, which treats design views, likewise takes the assignment of values to
each parameter as a presheaf and glues views that agree on the shared parameters. What is at
work here is the sheaf condition given by agreement on pairwise intersections
(\cite[\href{https://arxiv.org/html/2605.08609v1}{Definition 3.2, Theorem 4.1, \S\S3--5}]{Gibson26}).

Before gluing there is the question of whether the mismatches can be resolved.
The pattern of asking for the existence of a global extension also appears in the study of
contextuality. Abramsky and Brandenburger treat the distribution of outcomes for each
measurement context as a compatible empirical model, and put the extension to a global
distribution in correspondence with the existence of a factorizable hidden-variable model
(\cite[\href{https://arxiv.org/pdf/1102.0264v7}{\S\S2.2--2.5, Theorem 8.1}]{AB11}).

Abramsky, Mansfield, and Barbosa go further and linearize the support over a free abelian
group, constructing a \v{C}ech obstruction class for the extension of a local section.
A nonzero obstruction shows non-extendability. When the obstruction vanishes, on the other
hand, what one obtains is a compatible family after linearization, and an extension to an
actual assignment of outcomes requires an additional condition.
The Hardy model is an example that exhibits this difference
(\cite[\href{https://arxiv.org/pdf/1111.3620v2}{\S4, Propositions 4.2--4.3, \S5}]{AMB12}).

AAT constructs a global repair from the vanishing of the specified obstruction class, under
the torsor structure of repair states given by abelian-group coefficients, the equivariance
of restriction and action, and the sheaf condition (Theorem~\ref{thm:2.22}).
This is a construction that ties the standard fact---that a section of a torsor
trivializes it---to the evaluation of Laws and to repair states
(\cite[\href{https://stacks.math.columbia.edu/tag/03AG}{Tag 03AG}]{Stacks}).

The choice of coefficients also determines the meaning of a diagnosis.
Young's proposal for program analysis takes the substructures of a program as a site and
uses presheaves whose values are complete lattices of semantic properties
(\cite[\href{https://arxiv.org/html/2603.27015v1}{\S\S3.1--3.4}]{Young26}).

AAT compares abelian-group coefficients generated independently from the relations among
repair words and from the equations of the Laws.
We showed this comparison in the earlier work SAGA.
From an equivariant correspondence of local states we derive the soundness of the relations,
and under two completeness conditions we obtain a coefficient isomorphism and a
correspondence of cochains and residual classes. Adding the sheaf condition on repair states
leads to a judgment of global repair
(\cite[\href{https://arxiv.org/html/2608.21458v1}{\S\S3.4--3.7, 4--5, Theorems 5.1--5.2}]{SAGA}).

On the basis of this comparison from SAGA (Theorem~\ref{thm:2.49},
Corollary~\ref{cor:2.50}), AAT compares obstructions and diagnoses along a change of
reading. Under conditions on the relation between coefficients and on common
representatives, it reflects the vanishing of the original obstruction class from the
vanishing of the diagnosis (Theorem~\ref{thm:3.40}).
It further compares the two along a refinement of covers, and carries diagnostic invariance
back to the verdict on obstructions
(Theorems~\ref{thm:3.43} and~\ref{thm:3.44}).

\section{Abstraction and preservation of information}\label{sec:R.3}

The information that a readout must retain varies with the question one wants to answer.
In the abstract interpretation of Cousot and Cousot, the concrete and the abstract domains are connected by an order structure, and the correctness of a global fixpoint computation is
derived from approximations of local semantic operations. The relation between abstraction
and concretization underlies approximations that contain the actual behavior
(\cite[\href{https://www.di.ens.fr/~cousot/publications.www/CousotCousot-POPL-77-ACM-p238--252-1977.pdf}{\S\S5--7, pp. 240--243}]{CC77}).
Giacobazzi et al. study in more detail the form of completeness in which abstraction
commutes with the semantic operations, and distinguish it from completeness for the fixpoint result
(\cite[\href{https://www.sci.unich.it/~scozzari/paper/JACM00.pdf}{\S3, p. 372}]{GRS00}).
Besides what is approximated, what is preserved exactly for the specified operations becomes
the issue.

AAT constructs, as the canonical resolution, the coarsest reading that retains the values of
the Laws exactly (Theorem~\ref{thm:3.4}).
Preserving diagnoses requires, in addition, comparisons of coefficients and covers.
The isomorphism on first cohomology under Condition~C (Theorem~\ref{thm:3.17}) and the
reflection of the vanishing of the specified obstruction class (Theorem~\ref{thm:3.40})
capture the sufficiency of this information with different strengths.

AAT characterizes the information needed to judge a change by the kernel of the observation.
For a comparison $c$, let $\Gamma_c$ be the subgroup of changes that preserve it, and let
$O$ be the observation homomorphism. Membership in $\Gamma_c$ can be determined by
observations alone
if and only if $\ker O\subseteq\Gamma_c$ (Theorem~\ref{thm:7.9}).
If every change that observations do not distinguish preserves the comparison, then the same
observation yields the same judgment.

The sufficiency here concerns the judgment of a fixed property.
Local reconstruction, which recovers objects and all structure-preserving morphisms,
requires in addition the separation and assembly of local data (Chapter~\ref{chap:8}).

\section{Transport and base change}\label{sec:R.4}

If the schema changes, the data must be moved along with it.
Spivak's functorial data migration treats a schema as a category presented by a generating
graph and relations among paths, an instance as a set-valued functor, and a morphism of
instances as a natural transformation. Restriction along a functor of schemas has both a
left and a right adjoint, which yields three transfer functors
(\cite[\href{https://arxiv.org/pdf/1009.1166v3}{\S\S3.2, 3.4--3.5, 4, Propositions 4.2.1, 4.3.1}]{Spivak12}).
The algebraic databases of Schultz et al. incorporate types, operations, and equations into
this description, and construct instances that link entities with algebras on the type side,
together with restriction along a schema mapping and its left and right adjoints
(\cite[\href{https://arxiv.org/pdf/1602.03501v3}{\S\S6--7, especially \S6.19, Propositions 7.3--7.4}]{SSVW17}).

In this theory, schemas and their morphisms, instances, and queries fit into a common
double-categorical structure called a proarrow equipment. Transfers and the evaluation of
queries are also described using composition and adjunctions in that structure
(\cite[\href{https://arxiv.org/pdf/1602.03501v3}{Definition 8.13, Proposition 8.14, Lemma 8.18, \S\S8.25--8.26, 9}]{SSVW17}).
Treating objects and transfers in the same framework is a point of contact with AAT at
the level of the overall system.

AAT builds a reindexing from an exact change of reading, and constructs the transport and
the pullback of the relative core.
Its universal property rests on the standard theory of fibered categories
(\cite[\href{https://stacks.math.columbia.edu/tag/02XJ}{Tag 02XJ}]{Stacks}).
In contrast with the adjunctions of general data migration, it is this exactness condition
that makes the transport and the pullback of the core an adjoint equivalence
(Proposition~\ref{prop:5.8}).

Once two transfers are combined, a comparison of routes is needed in turn.
The method of forming a mate from the units and counits of adjunctions and of capturing its
invertibility as the Beck--Chevalley condition is made explicit in Shulman's study of framed
bicategories and fibrations
(\cite[\href{https://tac.mta.ca/tac/volumes/20/18/20-18.pdf}{\S13, Definition 13.11}]{Shulman08}).
AAT forms this mate for squares of readings generated from finite codes, and shows its
coherence with the comparison obtained by changing the choice of cartesian lift
(Theorem~\ref{thm:5.11}, Proposition~\ref{prop:5.12}).

From the square realized by finite codes and the common comparison diagram of
Construction~\ref{cons:5.37}, AAT builds the generated comparison of full geometries and
factors it into a composite of an invertible morphism and an idempotent.
Restricting to the image of the idempotent yields an invertible comparison
(Theorem~\ref{thm:6.16}).

\section{Lenses, updates, and comparison-preserving changes}\label{sec:R.5}

Even when the shipping address of an order is updated, we want the payment information to be
carried over unchanged. In updating a view, what matters is how to treat information that,
as here, does not appear in the readout.
Bancilhon and Spyratos specify that information as a complement, and define the translation
of an update by the condition that a fixed complement be preserved.
When a lift exists, that translation is unique.
For a family of updates that is closed under composition and in which updates can be undone
at each state, if all updates can be lifted under this condition, one obtains a translator
that preserves composition
(\cite[\href{https://www.inf.unibz.it/franconispace/lib/exe/fetch.php?media=organisation\%3Asakt\%3A2013\%3A5.pdf}{\S\S3, 5, 7, Theorems 5.6, 7.1}]{BS81}).

The lenses of Foster et al. treat bidirectional transformations through the laws of the read
operation get and the update operation put. The total lenses that AAT treats correspond to the very
well-behaved and total lenses of the original paper
(\cite[\href{https://www.cis.upenn.edu/~bcpierce/papers/newlenses-full.pdf}{\S3.1, Definitions 3.1.2--3.1.3, PutPut}]{FGMPS04}).
It is Johnson et al. who captured the relation between these laws and complements
category-theoretically. In a category with finite products, if the morphism from the view to
the terminal object splits, then the category of lenses is equivalent to the category of
complements. In the case of sets, once a reference point of the view is chosen, states can
be expressed as the product of the view and its reference fiber
(\cite[\href{https://mta.ca/~rrosebru/articles/Lens2Cambridge.pdf}{\S3, Propositions 3.1--3.2}]{JRW12}).

Using this product decomposition, AAT describes any morphism that preserves get and put as
the restriction and extension of a map of reference fibers
(Propositions~\ref{prop:1.33} and~\ref{prop:1.34}).
We connect this description to typed morphisms and to comparisons.

As for the universal property of lifting updates, Johnson et al. relate categorical lenses
to split opfibrations
(\cite[\href{https://mta.ca/~rrosebru/articles/Lens2Cambridge.pdf}{\S4, Remark 4.1, Corollary 4.1, Proposition 4.3}]{JRW12}).
AAT classifies the invertible changes that preserve a comparison under a fixed semantics.
For models in which the operations carry hidden states unchanged, AAT classifies the changes
of those states by the connected components of the operation graph, and expresses their
relation to visible changes by a split short exact sequence
(Theorems~\ref{thm:7.24} and~\ref{thm:7.25}).
For self-changes of the same product lens, under a fixed visible permutation the hidden
permutation agrees on all views.
For protocols that keep internal states, it can be chosen for each independent connected
component (Propositions~\ref{prop:7.27} and~\ref{prop:7.29}).
In these models, the way the operations tie states together determines the freedom of
change.

\section{Presentation and local reconstruction}\label{sec:R.6}

To treat objects and their changes by local presentations, structure-preserving morphisms
must be recoverable as well.
AAT assembles objects from local data and compatibility equations defined independently for
each input family, and puts local morphisms in correspondence with global
structure-preserving morphisms (Chapter~\ref{chap:8}).

At the basis of obtaining a category of models from operations and equations lies Lawvere's
functorial semantics. It treats an algebraic theory as a category with finite products, an
algebra as a product-preserving functor, and a homomorphism as a natural transformation
(\cite[\href{https://www.sas.rochester.edu/mth/sites/doug-ravenel/otherpapers/lawvere.pdf}{pp. 869--870}]{Lawvere63}).
The sketches of Barr and Wells give a graph together with diagrams to be made commutative
and specified cones. A model is an interpretation that makes the diagrams commute and sends
the cones to limit cones, and a morphism is a natural transformation
(\cite[\href{https://tac.mta.ca/tac/reprints/articles/12/tr12.pdf}{Chapter 4, \S1.2}]{BW85}).

The nerve theorem of Berger et al. constructs a fully faithful nerve from the category of
algebras to a presheaf category, under conditions on a dense generator and on a monad having
it as arities. Its essential image is characterized by a condition on restriction
(\cite[\href{https://math.univ-cotedazur.fr/~cberger/arities.pdf}{Theorem 1.10}]{BMW12}).

Here it also matters how presentations with the same meaning are identified.
The 2025 third version of \emph{Algebraic Databases} corrects the equivalence it had stated
between the category of presentations and the category of meanings, on the ground that
distinct syntactic morphisms can represent the same semantic morphism
(\cite[\href{https://arxiv.org/pdf/1602.03501v3}{Appendix B.1--B.3}]{SSVW17}).
AAT proves separately that a global morphism can be built from local maps and that morphisms
with the same readout coincide, and assembles local objects as well, up to isomorphism.
Deriving these three properties from the independent local conditions on each input family,
we obtain the equivalence between the category of realizations and the category of local
models. It is standard to derive an equivalence of categories from full faithfulness and
essential surjectivity
(\cite[\href{https://stacks.math.columbia.edu/tag/02C3}{Tag 02C3, Lemma 4.2.19}]{Stacks}).

Among the local conditions appear the graphs of get and put for lenses together with the
three laws, the graphs of generating edges and observations for protocols together with the
relations among paths, and the preservation conditions on coefficients and evaluation for
full geometries. For lenses we use finite reference fibers, and for protocols the finite
state set at each vertex and its finite list cover.

For the extension to images of idempotents we use the standard theory of Cauchy completion.
Kelly shows that, in the case enriched over sets, the small projectives of the presheaf
category over a small category are the retracts of representable presheaves, and that the
Cauchy completion corresponds to the completion that adds splittings of idempotents
(\cite[\href{https://www.sas.rochester.edu/mth/sites/doug-ravenel/otherpapers/kelly-book.pdf}{\S5.8, Theorem 5.36}]{Kelly82}).
AAT factors a generated comparison through the idempotent of normalization and builds the
restriction and extension to the image
(Theorems~\ref{thm:6.12}, \ref{thm:6.16}, and~\ref{thm:6.27}). The reconstruction of lenses
and protocols from reference fibers and generating tables also extends to the Karoubi
envelope and to the category of morphisms.

\section{Situating AAT}\label{sec:R.7}

The contribution of AAT lies in the theorems that connect its constructions and in the
correspondence with the semantics of individual problems.
The \textbf{Rising Sea} that AAT takes as its approach is to build foundations of this
kind and thereby to see individual problems as special cases of common constructions.
